\clearpage
\section{Typed boundary isomorphisms for renaming}
\label{sec:copying-boundary-iso}
\termdef{copying-boundary-iso}
The existing recursive graph isomorphism fixes each external boundary position.
Environment renaming may move those positions. We introduce a separate typed
boundary-isomorphism interface and prove a direct resource-permuting
automorphism of the original distributor. Neither the existing isomorphism nor
the observation relation is replaced. These are structural prerequisites for
whole-term renaming; the selected-resource clauses and propagation through
term constructors remain to be proved.

\subsection{The interface and its elementary laws}
A boundary isomorphism between two copying hypernets at the same allowance
consists of bijections on all vertices and on edges. The vertex bijection
preserves wire labels and preserves and reflects the predicates of being an
input and being an output. Hence it also preserves internal vertices. It may
permute external positions, but cannot exchange input and output roles or
mix term and atom wire types. Corresponding edge labels satisfy the existing
recursive label relation. Mapping an entire source or target list gives the
corresponding list in the target graph, in the same order. The definition
permits different interface lists when these conditions can be met.
\lean{Copying/BoundaryIso}{11}{19}{BoundaryIso}

Identity uses the identity bijections, reflexivity of label equivalence and
identity list maps. To reverse an isomorphism, use the two inverse bijections.
Wire labels and boundary roles are transported by their right-inverse laws.
Reverse the recursive edge-label relation. For an incidence list, apply the
inverse vertex map to its proved forward equality; cancellation of the two
vertex maps and the edge inverse identity give the required reversed equality.
\lean{Copying/BoundaryIso}{23}{31}{BoundaryIso.refl}
\lean{Copying/BoundaryIso}{33}{46}{BoundaryIso.symm}

Composition uses composition of the two vertex bijections and the two edge
bijections. Wire-label and boundary-role preservation compose directly;
recursive label equivalence is transitive. The list-map composition identity
reduces each incidence condition to the two component incidence conditions.
\lean{Copying/BoundaryIso}{48}{61}{BoundaryIso.trans}

Directed adjacency is preserved and reflected. A common vertex in a source
and target list maps to a common vertex in the image lists. Conversely,
recover preimages from the two mapped-list memberships; injectivity of the
vertex bijection identifies those preimages and supplies a common original
vertex. The statement therefore concerns actual incidence adjacency, not
just a relation between labels.
\lean{Copying/BoundaryIso}{64}{79}{BoundaryIso.adjacent_iff}

Every existing boundary-fixing recursive isomorphism embeds into the new
interface. Extend its internal-vertex bijection by the identity on input and
output constructors; reuse its edge bijection, label relation and incidence
proofs. For boundary ports, wire-label and role preservation are reflexive;
for internal ports use the original internal-label equation. This is an
embedding of the stronger boundary-fixing data, not a converse identification.
\lean{Copying/BoundaryIso}{81}{93}{BoundaryIso.ofDeepIso}

\subsection{An automorphism of the original distributor}
The original ordered parallel assembly provides a bijection between its
vertices and pairs consisting of a resource index and a component vertex.
Invert this bijection, apply the resource-tag permutation from Section~72,
and apply the assembly bijection again. This defines a bijection on all
vertices of the original distributor itself.
\lean{Copying/DistributorBoundaryIso}{9}{13}{Distributor.resourcePorts}

On a vertex embedded from component $i$, the inverse assembly cancels the
embedding. The result is exactly the same component vertex embedded at
$\pi(i)$. This cancellation identity proves the concrete boundary action:
the input for resource $i$ at use $j$ is sent to the input for $\pi(i)$ at the
same use $j$, and output $i$ is sent to output $\pi(i)$. Thus the input
positions transform as $i+kj\mapsto\pi(i)+kj$, matching the preceding
permuted-boundary construction.
\lean{Copying/DistributorBoundaryIso}{15}{23}{Distributor.resourcePorts_embed}
\lean{Copying/DistributorBoundaryIso}{26}{31}{Distributor.resourcePorts_input}
\lean{Copying/DistributorBoundaryIso}{33}{37}{Distributor.resourcePorts_output}

Use this all-vertex bijection and resource retagging on edges to construct
a boundary isomorphism from the original distributor to itself, at every
allowance. To prove wire-label preservation, express an arbitrary vertex
through the inverse assembly bijection and use the typing equation of each
component embedding. Input and output roles are preserved by a case split on
the component vertex. Edge labels are unchanged. For source and target lists,
compose the list maps and apply the embedding cancellation identity entry
by entry. All canonical weakening and contraction edges remain, with their
original ordered incidences; no positivity assumption on the resource or
use count is needed.
\lean{Copying/DistributorBoundaryIso}{41}{83}{Distributor.boundaryIso}

This interface makes the external action explicit on the original graph,
whereas Section~72 compared a separately boundary-permuted assembly with the
original assembly using boundary-fixing isomorphism. Neither result proves
observational equivalence under arbitrary boundary permutations. The original
terms and translation of Proposition~12.2 remain unchanged, and whole-term
renaming, substitution and the copying machine remain outstanding.
